\chapter{Introducción}
\section{Motivación y contexto}
En las últimas décadas, el incremento de la esperanza de vida ha traído consigo un aumento significativo en la prevalencia de enfermedades neurodegenerativas. Entre ellas, la enfermedad de Alzheimer constituye la principal causa de demencia a nivel global, afectando no solo a quienes la padecen, sino también a su entorno familiar y social. Según la Organización Mundial de la Salud, se estima que más de 55 millones de personas viven con demencia en el mundo, y esta cifra continúa en ascenso debido al envejecimiento poblacional \cite{who2021dementia}.

El Alzheimer se caracteriza por un deterioro progresivo de las funciones cognitivas, comenzando generalmente con pérdida de memoria a corto plazo y evolucionando hacia la pérdida total de autonomía \cite{merino2015enfermedad}. Aunque la investigación médica ha logrado avances significativos en la comprensión de los mecanismos neuropatológicos de la enfermedad, en la actualidad no existe un tratamiento curativo. Las intervenciones disponibles se centran en ralentizar el progreso de los síntomas y mejorar la calidad de vida de los pacientes, combinando tratamiento farmacológico con terapias no farmacológicas como la estimulación cognitiva \cite{merino2015enfermedad}.

La estimulación cognitiva asistida por ordenador ha demostrado ser una herramienta eficaz para preservar las capacidades intelectuales y retrasar el deterioro funcional \cite{kueider2012computerized}. Diversos estudios avalan que la práctica regular de ejercicios cognitivos estructurados favorece la neuroplasticidad y la sinaptogénesis, contribuyendo a mantener activas las redes neuronales \cite{castillo2020neuroplasticidad}. Sin embargo, las soluciones comerciales actuales presentan limitaciones significativas: elevados costes de suscripción, falta de adaptación al idioma y contexto cultural español, y una orientación predominante hacia dispositivos móviles que no siempre resulta accesible para el perfil de usuario mayor \cite{dort2020_aphasia, stroke_equity_2021}.

En este contexto surge ReMind, una aplicación web diseñada específicamente para ofrecer estimulación cognitiva a personas mayores, con especial atención a aquellas en fases tempranas de Alzheimer o en riesgo de desarrollarlo. La plataforma se concibe como una herramienta gratuita, accesible desde ordenadores de sobremesa, que integra minijuegos terapéuticos basados en ejercicios validados científicamente. El sistema permite además la supervisión por parte de terapeutas, quienes pueden gestionar agendas personalizadas de actividades y realizar un seguimiento objetivo del progreso de cada paciente.

\section{Objetivos}
\subsection{Objetivo general}
El objetivo principal de este proyecto es diseñar, implementar y validar una aplicación web que proporcione un conjunto de juegos y actividades de estimulación cognitiva orientados a prevenir o ralentizar el deterioro asociado al Alzheimer. La plataforma debe ofrecer una experiencia accesible e intuitiva, adaptada a las necesidades de un público mayor, e integrar funcionalidades de gestión clínica que permitan a los terapeutas monitorizar la evolución de sus pacientes.

\subsection{Objetivos específicos}
Para alcanzar el objetivo general, se definen los siguientes objetivos específicos:
\begin{itemize}
    \item Investigar los fundamentos neurocientíficos de la enfermedad de Alzheimer, incluyendo sus síntomas, causas y las líneas de tratamiento actuales, con énfasis en las terapias de estimulación cognitiva.
    \item Estudiar las distintas áreas del desarrollo cognitivo y las técnicas de estimulación validadas para identificar los tipos de ejercicios más adecuados para su implementación informática.
    \item Diseñar una arquitectura software robusta, escalable y segura, siguiendo el patrón cliente-servidor de tres capas, que garantice la separación de responsabilidades y la mantenibilidad del sistema.
    \item Desarrollar un conjunto de minijuegos interactivos que trabajen diferentes dominios cognitivos como la memoria, la atención, el lenguaje, el cálculo y las funciones ejecutivas.
    \item Implementar un sistema de gestión de usuarios con tres roles diferenciados (administrador, terapeuta y paciente) que permita la asignación personalizada de actividades y el control de acceso basado en privilegios.
    \item Diseñar una interfaz de usuario accesible que cumpla con los estándares WCAG 2.1, priorizando la usabilidad para personas mayores con potenciales déficits visuales, motrices o cognitivos.
    \item Incorporar un módulo de seguimiento y estadísticas que permita a los terapeutas evaluar el rendimiento y la evolución de los pacientes a lo largo del tiempo.
\end{itemize}

